\documentclass[10pt,a4paper]{amsart}
\usepackage{amsmath,amssymb,amsthm,graphicx}
\usepackage{verbatim}
\newenvironment{reference}{\comment}{\endcomment}
%\newenvironment{reference}{}{}
\newenvironment{slogan}{\comment}{\endcomment}
\newenvironment{history}{\comment}{\endcomment}

% For commutative diagrams we use Xy-pic
\usepackage[all]{xy}

% We use 2cell for 2-commutative diagrams.
\xyoption{2cell}
% Theorem environments.
%
\theoremstyle{plain}
\newtheorem{theorem}[subsection]{Theorem}
\newtheorem{proposition}[subsection]{Proposition}
\newtheorem{lemma}[subsection]{Lemma}

\theoremstyle{definition}
\newtheorem{definition}[subsection]{Definition}
\newtheorem{example}[subsection]{Example}
\newtheorem{exercise}[subsection]{Exercise}
\newtheorem{situation}[subsection]{Situation}

\theoremstyle{remark}
\newtheorem{remark}[subsection]{Remark}
\newtheorem{remarks}[subsection]{Remarks}

\numberwithin{equation}{subsection}

% Macros
%
\def\lim{\mathop{\mathrm{lim}}\nolimits}
\def\colim{\mathop{\mathrm{colim}}\nolimits}
\def\Spec{\mathop{\mathrm{Spec}}}
\def\Hom{\mathop{\mathrm{Hom}}\nolimits}
\def\Ext{\mathop{\mathrm{Ext}}\nolimits}
\def\SheafHom{\mathop{\mathcal{H}\!\mathit{om}}\nolimits}
\def\SheafExt{\mathop{\mathcal{E}\!\mathit{xt}}\nolimits}
\def\Sch{\mathit{Sch}}
\def\Mor{\mathop{\mathrm{Mor}}\nolimits}
\def\Ob{\mathop{\mathrm{Ob}}\nolimits}
\def\Sh{\mathop{\mathit{Sh}}\nolimits}
\def\NL{\mathop{N\!L}\nolimits}
\def\CH{\mathop{\mathrm{CH}}\nolimits}
\def\proetale{{pro\text{-}\acute{e}tale}}
\def\etale{{\acute{e}tale}}
\def\QCoh{\mathit{QCoh}}
\def\Ker{\mathop{\mathrm{Ker}}}
\def\Im{\mathop{\mathrm{Im}}}
\def\Coker{\mathop{\mathrm{Coker}}}
\def\Coim{\mathop{\mathrm{Coim}}}

% Boxtimes
%
\DeclareMathSymbol{\boxtimes}{\mathbin}{AMSa}{"02}

%
% Macros for moduli stacks/spaces
%
\def\QCohstack{\mathcal{QC}\!\mathit{oh}}
\def\Cohstack{\mathcal{C}\!\mathit{oh}}
\def\Spacesstack{\mathcal{S}\!\mathit{paces}}
\def\Quotfunctor{\mathrm{Quot}}
\def\Hilbfunctor{\mathrm{Hilb}}
\def\Curvesstack{\mathcal{C}\!\mathit{urves}}
\def\Polarizedstack{\mathcal{P}\!\mathit{olarized}}
\def\Complexesstack{\mathcal{C}\!\mathit{omplexes}}
% \Pic is the operator that assigns to X its picard group, usage \Pic(X)
% \Picardstack_{X/B} denotes the Picard stack of X over B
% \Picardfunctor_{X/B} denotes the Picard functor of X over B
\def\Pic{\mathop{\mathrm{Pic}}\nolimits}
\def\Picardstack{\mathcal{P}\!\mathit{ic}}
\def\Picardfunctor{\mathrm{Pic}}
\def\Deformationcategory{\mathcal{D}\!\mathit{ef}}
\usepackage{cleveref}
\crefname{originalreference}{original reference}{original references}
\Crefname{originalreference}{Original reference}{Original references}
\makeatletter
\newlabel{lemma-exact-sequence-j-star}{{08BV}{source}{}{sourcepointer-1}{}}
\newlabel{lemma-exact-sequence-j-star@cref}{{[originalreference][0][]08BV}{[1][0][]source}{}{}{}}
\makeatother
\crefname{originalreference}{original reference}{original references}
\Crefname{originalreference}{Original reference}{Original references}
\makeatletter
\newlabel{lemma-exact-sequence-lower-shriek}{{08BU}{source}{}{sourcepointer-2}{}}
\newlabel{lemma-exact-sequence-lower-shriek@cref}{{[originalreference][0][]08BU}{[1][0][]source}{}{}{}}
\makeatother
\crefname{originalreference}{original reference}{original references}
\Crefname{originalreference}{Original reference}{Original references}
\makeatletter
\newlabel{lemma-make-more-elementary-distinguished-squares}{{08GN}{source}{}{sourcepointer-3}{}}
\newlabel{lemma-make-more-elementary-distinguished-squares@cref}{{[originalreference][0][]08GN}{[1][0][]source}{}{}{}}
\makeatother
\crefname{originalreference}{original reference}{original references}
\Crefname{originalreference}{Original reference}{Original references}
\makeatletter
\newlabel{lemma-restrict-direct-image-open}{{08FE}{source}{}{sourcepointer-4}{}}
\newlabel{lemma-restrict-direct-image-open@cref}{{[originalreference][0][]08FE}{[1][0][]source}{}{}{}}
\makeatother
\crefname{originalreference}{original reference}{original references}
\Crefname{originalreference}{Original reference}{Original references}
\makeatletter
\newlabel{lemma-restrict-lower-shriek}{{08HA}{source}{}{sourcepointer-5}{}}
\newlabel{lemma-restrict-lower-shriek@cref}{{[originalreference][0][]08HA}{[1][0][]source}{}{}{}}
\makeatother
\crefname{originaltheorem}{original theorem}{original theorems}
\Crefname{originaltheorem}{Original theorem}{Original theorems}
\makeatletter
\newlabel{spaces-properties-theorem-exactness-stalks}{{04K5}{source}{}{sourcepointer-6}{}}
\newlabel{spaces-properties-theorem-exactness-stalks@cref}{{[originaltheorem][0][]04K5}{[1][0][]source}{}{}{}}
\makeatother
\title{Pushforward and shriek maps are isomorphisms for objects supported off the distinguished-square legs}
\author{The Stacks Project Authors}
\date{stacks-project:spaces-perfect.tex:local-snapshot-20261012; GNU FDL 1.2 or later}
\begin{document}
\maketitle
\noindent\textit{Editorial note.} All delivered mathematical text is the author's own, copied from the pinned Stacks Project chapter file spaces-perfect.tex; no statement or proof step was rewritten, reordered or removed.\par\medskip
\noindent\textit{Editorial note.} The counted proof is the author's complete argument: a direct part (1) from exact sequences, support statements for both pushforwards, a section-reduction of the stalk verification, and an explicit section-case computation, all inside the delivered span\par\medskip
\noindent\textit{Editorial note.} The label lemma-pushforward-with-support-in-open names the counted statement, and the proof environment immediately following it in the source is delivered verbatim as its proof.\par\medskip
\noindent\textit{Editorial note.} The counted statement is carried by the author's own lemma environment in the chapter source.\par\medskip
\noindent\textit{Allowed standard prerequisites (editorial).}\begin{itemize}
\item Elementary distinguished squares and the exact sequences relating j\_!, Rj\_* and restriction to the two legs
\item Restriction formulae for Rj\_* and j\_! over the legs of a distinguished square
\item The stalk criterion for isomorphisms in the derived category of algebraic spaces, and base change of distinguished squares
\end{itemize}\medskip
\ifdefined\hypertarget\hypertarget{sourcepointer-1}{}\fi
\noindent\textit{Source reference.} Original reference 08BV: sources/stacks-spaces-perfect/pinned-original-source.tex, source lines 1889--1889. Prerequisite pointer invoked at proof line(s) 2152; it names a standard fact of the same chapter that the counted argument uses, and every step of the counted proof built on it is retained verbatim in the delivered body.\par\medskip
\ifdefined\hypertarget\hypertarget{sourcepointer-2}{}\fi
\noindent\textit{Source reference.} Original reference 08BU: sources/stacks-spaces-perfect/pinned-original-source.tex, source lines 1823--1823. Prerequisite pointer invoked at proof line(s) 2154; it names a standard fact of the same chapter that the counted argument uses, and every step of the counted proof built on it is retained verbatim in the delivered body.\par\medskip
\ifdefined\hypertarget\hypertarget{sourcepointer-3}{}\fi
\noindent\textit{Source reference.} Original reference 08GN: sources/stacks-spaces-perfect/pinned-original-source.tex, source lines 1541--1541. Prerequisite pointer invoked at proof line(s) 2176; it names a standard fact of the same chapter that the counted argument uses, and every step of the counted proof built on it is retained verbatim in the delivered body.\par\medskip
\ifdefined\hypertarget\hypertarget{sourcepointer-4}{}\fi
\noindent\textit{Source reference.} Original reference 08FE: sources/stacks-spaces-perfect/pinned-original-source.tex, source lines 275--275. Prerequisite pointer invoked at proof line(s) 2159; it names a standard fact of the same chapter that the counted argument uses, and every step of the counted proof built on it is retained verbatim in the delivered body.\par\medskip
\ifdefined\hypertarget\hypertarget{sourcepointer-5}{}\fi
\noindent\textit{Source reference.} Original reference 08HA: sources/stacks-spaces-perfect/pinned-original-source.tex, source lines 2119--2119. Prerequisite pointer invoked at proof line(s) 2162; it names a standard fact of the same chapter that the counted argument uses, and every step of the counted proof built on it is retained verbatim in the delivered body.\par\medskip
\ifdefined\hypertarget\hypertarget{sourcepointer-6}{}\fi
\noindent\textit{Source reference.} Original theorem 04K5: sources/stacks-chow/tags.txt, source lines 5961--5961. Prerequisite pointer invoked at proof line(s) 2170; it names a standard fact of a sibling chapter that the counted argument uses, and every step of the counted proof built on it is retained verbatim in the delivered body.\par\medskip
\begin{lemma}
\label{lemma-pushforward-with-support-in-open}
Let $S$ be a scheme. Let $(U \subset X, j : V \to X)$ be an elementary
distinguished square of algebraic spaces over $S$. Set
$T = |X| \setminus |U|$.
\begin{enumerate}
\item If $E$ is an object of $D(\mathcal{O}_X)$ supported on $T$, then
(a) $E \to Rj_*(E|_V)$ and (b) $j_!(E|_V) \to E$ are isomorphisms.
\item If $F$ is an object of $D(\mathcal{O}_V)$ supported on $j^{-1}T$, then
(a) $F \to (j_!F)|_V$, (b) $(Rj_*F)|_V \to F$, and (c)
$j_!F \to Rj_*F$ are isomorphisms.
\end{enumerate}
\end{lemma}

\begin{proof}
Let $E$ be an object of $D(\mathcal{O}_X)$ whose cohomology sheaves are
supported on $T$. Then we see that $E|_U = 0$ and $E|_{U \times_X V} = 0$
as $T$ doesn't meet $U$ and $j^{-1}T$ doesn't meet $U \times_X V$.
Thus (1)(a) follows from Lemma \ref{lemma-exact-sequence-j-star}.
In exactly the same way (1)(b) follows from
Lemma \ref{lemma-exact-sequence-lower-shriek}.

\medskip\noindent
Let $F$ be an object of $D(\mathcal{O}_V)$ whose cohomology sheaves
are supported on $j^{-1}T$. By
Lemma \ref{lemma-restrict-direct-image-open} we have
$(Rj_*F)|_U = Rj_{W, *}(F|_W) = 0$ because $F|_W = 0$ by our assumption.
Similarly $(j_!F)|_U = j_{W!}(F|_W) = 0$ by
Lemma \ref{lemma-restrict-lower-shriek}.
Thus $j_!F$ and $Rj_*F$ are
supported on $T$ and $(j_!F)|_V$ and $(Rj_*F)|_V$ are supported on
$j^{-1}(T)$. To check that the maps (2)(a), (b), (c) are isomorphisms
in the derived category, it suffices to check that these map induce
isomorphisms on stalks of cohomology sheaves at geometric points of $T$
and $j^{-1}(T)$ by
Properties of Spaces, Theorem
\ref{spaces-properties-theorem-exactness-stalks}.
This we may do after replacing $X$ by $V$, $U$ by $U \times_X V$,
$V$ by $V \times_X V$ and $F$ by $F|_{V \times_X V}$ (restriction via
first projection), see
Lemmas \ref{lemma-restrict-direct-image-open},
\ref{lemma-restrict-lower-shriek}, and
\ref{lemma-make-more-elementary-distinguished-squares}.
Since $V \times_X V \to V$ has a section this
reduces (2) to the case that $j : V \to X$ has a section.

\medskip\noindent
Assume $j$ has a section $\sigma : X \to V$.
Set $V' = \sigma(X)$. This is an open subspace of $V$.
Set $U' = j^{-1}(U)$. This is another open subspace of $V$.
Then $(U' \subset V, V' \to V)$ is an elementary distinguished
square. Observe that $F|_{U'} = 0$ and $F|_{V' \cap U'} = 0$
because $F$ is supported on $j^{-1}(T)$. Denote $j' : V' \to V$
the open immersion and $j_{V'} : V' \to X$ the composition
$V' \to V \to X$ which is the inverse of $\sigma$.
Set $F' = \sigma^*F$. The distinguished triangles of
Lemmas \ref{lemma-exact-sequence-lower-shriek} and
\ref{lemma-exact-sequence-j-star} show that
$F = j'_!(F|_{V'})$ and $F = Rj'_*(F|_{V'})$.
It follows that $j_!F = j_!j'_!(F|_{V'}) = j_{V'!}F = F'$
because $j_{V'} : V' \to X$ is an isomorphism and the inverse
of $\sigma$. Similarly, $Rj_*F = Rj_*Rj'_*F = Rj_{V', *}F = F'$.
This proves (2)(c). To prove (2)(a) and (2)(b) it suffices
to show that $F = F'|_V$. This is clear because both $F$ and $F'|_V$
restrict to zero on $U'$ and $U' \cap V'$ and the same object
on $V'$.
\end{proof}

\end{document}
